\documentclass[11pt,a4paper]{article}
\usepackage[top=42pt,bottom=42pt,left=70pt,right=70pt]{geometry}  % to make pages wider

\usepackage{amsmath}
\usepackage{amsfonts}
\usepackage{graphicx}  % for figures etc.
\usepackage{url}  % to handle URLs 
\usepackage[comma,authoryear]{natbib}  % for author-date ref style --  allows \citet (textual), and \citep (parenthesized)
\usepackage[hidelinks]{hyperref}
\hypersetup{allcolors=[rgb]{0, 0, 0.33},colorlinks=true}

\renewcommand {\cite} {\citep}  % default for cite is citet in natbib - changes it to citep
\renewcommand\bibname{References}    % if not using newrucsthesis sty file

\usepackage{booktabs}
\usepackage{array}
\usepackage{enumitem}
\usepackage{amssymb}

\usepackage{setspace}
\doublespacing

\title{Methodology}
\author{g23n7880 }
\date{September 2026}

\begin{document}

\maketitle

\subsection{Dataset and Data Selection}
\label{sec:data_selection}

The construction of a supervised orthographic classifier requires training data in
which the orthographic variant of each sentence is known with confidence. For a
low-resource language such as Sesotho, no pre-existing labelled dataset covering
both the South African Sesotho (SAS) and Lesotho Sesotho (LS) orthographic
conventions simultaneously was available at the time of this study. This section
documents the sources selected to constitute the training corpus, the rationale
for each selection decision, and the labelling strategy employed to assign
orthographic variant labels without manual sentence-level annotation.

\subsubsection{Labelling strategy: source-based provenance}

Ground truth labels for classifier training were assigned using a
\textit{source-based provenance} strategy: sentences inherit the orthographic
variant label of the institutional source from which they were collected, rather
than from direct manual inspection of each sentence. [read MINZ 2009 paper]

This condition was verified empirically prior to classifier training. For each
source, the average occurrence of variant-discriminating graphemic markers
per sentence was computed using the substitution rules documented by
\citet{makutoane2022people}. Specifically, the semi-vowel pairs
\textit{ea}/\textit{oa} (LS) versus \textit{ya}/\textit{wa} (SAS), and
the consonant cluster pairs \textit{kh}/\textit{ch} (LS) versus
\textit{kg}/\textit{tjh} (SAS). LS-labelled sentences yielded a mean LS
marker density of 0.89 per sentence and a mean SAS marker density of 0.00.
SAS-labelled sentences yielded a mean SAS marker density of 1.34 per sentence
and a mean LS marker density of 0.00. The complete absence of cross-variant
markers in each group confirms that source provenance reliably identifies
orthographic variant and that source-based labelling is a valid ground truth
strategy for this corpus.

\subsubsection{South African Sesotho sources}

\paragraph{Vuk'uzenzele — DSFSI corpus.}
The primary source of SAS training data is the Vuk'uzenzele corpus maintained
by the Data Science for Social Impact (DSFSI) research group at the University
of Pretoria \citep{marivate2020investigating}. Vuk'uzenzele is the South African
government's official public information magazine, published in all eleven
official languages including Sesotho. The magazine is produced under South
African government style guidelines, which mandate adherence to the SAS
orthographic standard formalised in 1959. The DSFSI corpus provides
sentence-aligned English--Sesotho pairs extracted from digitised magazine
editions spanning 2018 to 2022, representing 84 editions. The Sesotho component
of the simple sentence-alignment output was used, extracting the target-language
column from each aligned CSV file.

This source was selected on three grounds. First, its institutional origin
guarantees SAS orthographic consistency, since government publications are
produced to a standardised style. Second, it represents the domain of public
information and civic communication which is underrepresented in existing
Sesotho resources, which are heavily skewed toward religious and legislative
text \citep{kwon2026afroscope}. Third, the corpus is publicly available under an
open licence, ensuring reproducibility.

\subsubsection{Lesotho Sesotho sources}

\paragraph{Sesotho News (SN) dataset.}
The primary source of LS training data is the Sesotho News dataset introduced
by \citet{mokhosi2024sesotho}, available from Zenodo (record 10531959). The
dataset comprises 2,401 news headlines manually scraped from two Lesotho-based
online newspapers and annotated for sentiment analysis and aspect-based sentiment
analysis. Headlines were collected using the Lesotho orthographic standard, as
confirmed by the source publication. Two files were used from the dataset:
\texttt{NewsSA.txt} and \texttt{NewsABSA.txt}, providing sentence-level LS text
across a news domain.

This source was selected because it is the only existing Sesotho dataset that
explicitly documents its orthographic standard, is independently peer-reviewed,
covers the news domain, and provides sufficient sentence-level granularity for
classifier training. Its use here is consistent with its intended purpose: the
authors explicitly release it for NLP research on Sesotho. Notably, the same
dataset is used later as the evaluation benchmark for the downstream sentiment
classification task, creating a direct methodological connection between the
training corpus and the extrinsic quality evaluation.

\paragraph{Nalibali children's stories.}
A supplementary LS source was obtained by scraping the Nalibali reading
campaign website (\url{https://nalibali.org}), a South African literacy
initiative that publishes multilingual children's stories in all eleven official
languages. The Sesotho stories are available at the \texttt{?lang=st-ls} URL
parameter, denoting the Lesotho Sesotho orthographic variant. Story text was
extracted from the \texttt{elementor-widget-text-editor} content containers
using BeautifulSoup, with navigation elements, author attributions, and
discussion questions excluded via pattern matching.

This source was selected to increase LS domain diversity. The SN dataset
provides news headlines, which are short, syntactically compressed, and
domain-specific. Children's stories provide longer, more syntactically complete
sentences across narrative and descriptive registers, partially mitigating the
domain skew that \citet{kwon2026afroscope} identify as a systemic quality
failure in African language corpora. Nalibali is a registered non-governmental
organisation with a citable institutional identity, satisfying the requirement
for source-based provenance reliability.

\subsubsection{Corpus composition}

Following GlotLID verification, deduplication, and a minimum length filter of
five words, the training corpus comprised 302 sentences: 200 SAS sentences
from Vuk'uzenzele and 102 LS sentences from the SN dataset. The class
imbalance of approximately 2:1 reflects the relative availability of
digitised text in each orthographic variant rather than a sampling decision,
and is addressed during classifier training through balanced class weighting
rather than oversampling, to avoid introducing artificial distributional
regularities into the feature space. The corpus size, while modest relative
to high-resource standards, is appropriate for a proof-of-concept classifier
demonstrating the viability of the orthographic distinction detection approach
on a task where, as Section 3.4 establishes, the discriminating signal is
concentrated in a small and consistent set of sub-word character patterns. 

The corpus has grown to around 5000 sentence


\section{Implementation}
The pipeline was implemented in Python 3.11 across three modular scripts, each corresponding to a discrete pipeline stage. The modular architecture follows design principle articulated by \citet{suchomel2020better} that pipeline components should be independently testable before integration, so that errors introduced at one stage do not propagate silently to another stage. A similar modular approach is adopted by \citet{de2024data} in the Labadian Crawler for Tetun, and by \citet{gumede2025creating} in the ULPDO isiZulu corpus construction, both of which combine automated processing with targeted human input at key quality decision points.

\subsection{Stage 1: Semi-Automated Collection and Identification}
Data collection followed a semi-automaated approach in which  human judgement was exercised at the source-selection stage, with all subsequent processing automated. This design is consistent with \citet{gumede2025creating} hybrid approach as humans provide expert opinion on which sources are appropriate while automation handles extraction, verification, and structuring at scale. The distinction was made due to a narrowing of the scope for the project. The semi automated process is favoured by \citet{barbaresi2013challenges} as \citet{barbaresi2013challenges} argues that crawling without expert knowledge of the target language is liable to produce systematic distortions, since search engines and broad crawlers favour English content and structurally well-connected domains. For a language like Sesotho, whose digital presence is concentrated in specific institutional and community domains rather than broadly distributed across the web, targeted source selection by a research with knowledge of the language community provides more a reliable starting point than unrestricted automated discovery.

Three seed sources were selected, as described in Section~\ref{sec: data_selection}. For each source, text was retrieve using the Python \texttt{requests} library with a standard browser user-agent string to avoid bot rejection. HTML content was parsed using BeautifulSoup4, with content extracted from semantically appropriate containers identified through prior page inspection for each source. For the NaliBali source, which renders content through an Elementor page builder framework, text was extracted from \texttt{div} elements carrying the CSS class \texttt{elementor-widget-text-editor}, identified through diagnostic inspection of the page structure. 

Sentence segmentation was performed by splitting on terminal punctuation characters which are full stop, exclamation mark, question mark, and newline. Segements of more than three words or more were retained which excluded isolated headers, navigation lables, single-word or two-word fragments that are not linguistically usable while retaining grammatically complete short sentences typical of new headline and children's story text. \citep{de2024data} adopt a comparable minimu word threshold in their Tetun corpus pipeline.

Each candidate sentence was submitted individually to GlotLID \citep{kargaran2024glotcc}, a fastText-based language identification model covering 1655 language varieties. GlotLID was selected for its false positive rate minimisation over raw F1 score which is critical because incorrectly included text produces irreversible contamination while incorrectly excluded text is a recoverable loss \citep{kreutzer-etal-2022-quality}. This asymmetry is particularly consequential for low resource languages like Sesotho, which \citet{kargaran2024glotcc} demonstrate is prone to cousin-language misclassifications when LID tool with inadequate coverage are used. Sentences where GlotLID's top-1 prediction was \texttt{sot\_Latn} with a confidence score at or above 0.7 were retained; all others were discarded.

The confidence threshold of 0.7 was determined empirically through diagnostic testing on known Sesotho sentences from the Mokhosi dataset prior to full collection. Short headlines of four to six words recieved GlotLID scores between 0.85 and 0.91, indicating that the conventional threshold of 0.9, calibrated against longer texts in high-resource settings, would systematically reject valid LS content. \citet{fedorova2026openlid} document an equivalent precision-coverage trade-off in their evaluation of OpenLID-v3, finding that improving precision threshold removed up to 45\% of samples for minority languga varieties. The 0.7 threshold preserves short headline sentences while maintaining reliable separation from closely related languages. Diagnostic  evaluation confirmed that Sestswana (\texttt{tesn\_Latn}) and Sesotho sa Leboa (\texttt{nso\_Latn}) consistently scored below 0.2 on Sesotho-labelled sentences, providing a clear confidence margin between acceptance and rejection.
\subsection{Stage 2: Cleaning}
The cleaning stage applies four operations in sequence, each
motivated by a specific quality concern identified in the literature.

\textbf{Unicode normalisation to NFC form} was applied first using
Python's \texttt{unicodedata.normalize} function. NFC normalisation
ensures that diacritical marks are represented in their canonical
composed form rather than as decomposed sequences of base character
and combining diacritic. This step directly addresses the sub-lexical
degradation documented by \citet{de2007automatic}, who
demonstrate that OCR processing and non-standard text encoding
pipelines systematically strip or misrepresent diacritical marks in
African language text. For Lesotho Sesotho specifically, diacritical
preservation is not merely a formatting concern: LS orthography
encodes phonological distinctions through diacritics that SAS
orthography omits \citep{matlosa2017sesotho}, meaning that incorrect
normalisation would destroy a primary signal for the orthographic
classifier downstream.

\textbf{URL removal} was applied using a regular expression matching
\texttt{http} or \texttt{https} strings, eliminating hyperlink
artefacts introduced during web extraction that would otherwise
contribute meaningless character sequences to the feature space.

\textbf{Whitespace normalisation} collapsed multiple consecutive
spaces, tabs, and newlines into single spaces, and stripped leading
and trailing whitespace. This step addresses formatting artefacts
introduced by HTML rendering and line-break
conventions that vary across source formats.

\textbf{Deduplication} removed exact duplicate sentences using
pandas \texttt{drop\_duplicates} on the text column.
\citet{tatariya2024good} document that up to 40\% of
web-crawled content consists of exact duplicates, with
paragraph-level duplication reaching 70\% in some Common Crawl
snapshots. Deduplication at the sentence level ensures that
repetitive boilerplate content such as navigation text, repeated
disclaimers, and templated phrases common in government and
institutional sources does not inflate corpus statistics or
bias classifier training by overrepresenting particular character
sequences. This concern is also documented by
\citet{cahyawijaya2023nusawrites}, who observe that Buginese
Wikipedia data exhibited very low lexical diversity despite a large
token count, attributable to boilerplate repetition that deduplication
would have addressed.

A \textbf{minimum length filter} of five words was then applied,
removing fragments that do not constitute usable linguistic content.
The five-word threshold is consistent with the minimum meaningful
sentence length assumed in low-resource corpus construction
pipelines \citep{de2024data}.

\subsection{Stage 3: Orthographic Classification}
The orthographic classification stage consitutes the novel contribution of the piepline where no tool automatically distinguishes SAS from LS orthography at the sentence level. This stage was implmented as a two-component system being a rule-based basline classifer and a trained character n-gram classifier, with both evaluated on a held-out test set to assess the value added by the learned model over the explicit rule encoding.

\subsubsection{Rule-Based Baseline}
\label{sec:rulebased}

The rule-based classifier implements the substitution rules documented by \citet{makutoane2022people} directly as Python
marker-counting logic. For each sentence, the occurrence count of LS
markers (\textit{ea}, \textit{oa}, \textit{kh}, \textit{ch},
\textit{li}) and SAS markers (\textit{ya}, \textit{wa},
\textit{kg}, \textit{tjh}) are computed, and the sentence is
assigned the variant label corresponding to the higher marker count,
with SAS assigned as the default in the case of a tie. The marker
sets were derived directly from Makutoane's documented substitution
table, requiring no training data.

This baseline serves two purposes. First, it operationalises the
existing linguistic knowledge about the SAS/LS distinction,
providing a performance ceiling achievable through explicit rule
encoding. Second, it provides a principled lower bound against which
the learned classifier can be evaluated. If the trained model does
not surpass the rule-based baseline, this would indicate that the
character n-gram features capture nothing beyond what the literature
already documents.

\subsubsection{Character N-Gram Logistic Regression Classifier}
\label{sec:learned}

The learned classifier was implemented as a scikit-learn
\texttt{Pipeline} combining a \texttt{CountVectorizer} with a
\texttt{LogisticRegression} estimator. The pipeline abstraction
ensures that the same vectorisation parameters applied during training
are applied identically during inference on new text, eliminating
the risk of train-test leakage through inconsistent preprocessing.

\paragraph{Feature representation.}
Character n-grams spanning lengths two to four
(\texttt{ngram\_range=(2,4)}, \texttt{analyzer='char\_wb'}) were
selected as the feature representation. The character word-boundary
analyser (\texttt{char\_wb}) pads each token with spaces before
extracting n-grams, ensuring that boundary characters are captured
as informative features rather than discarded.

Character-level features are the appropriate representation for
the SAS/LS classification task because the orthographic distinction
is sub-lexical. The discriminating signal lies in character sequences
such as \textit{ea} and \textit{oa} (LS) versus \textit{ya} and
\textit{wa} (SAS), which appear within words rather than as
independent tokens. Word-level features would require the exact
surface form of a discriminating word to appear in a sentence before
any signal is captured; character n-grams capture orthographic
patterns regardless of the specific word in which they appear.
This property is critical for a morphologically rich agglutinative
language like Sesotho, where a single root generates thousands of
inflected surface forms through prefixation and suffixation
\citep{kambarami2021computational}. Character n-gram models have
established precedent in the similar-language and dialect
identification literature: \citet{cianflone2016n} report that
character n-gram models achieved 88.45\% accuracy on the VarDial
2016 Discriminating Similar Languages shared task, outperforming a
character-based convolutional neural network on the same data and
closely approaching the top system at 89.38\%. Their analysis
confirms that n-grams of size 7--8 peak in accuracy for similar
language discrimination, though for a binary classification task
with highly consistent orthographic signal, shorter n-grams of
size 2--4 are sufficient and reduce sparsity.
\citet{de2024data} similarly demonstrate that a Multinomial
Naive Bayes classifier with 5-gram characters achieved 99.77\%
accuracy for Tetun language identification, confirming character
n-gram effectiveness for low-resource language discrimination.

The n-gram range of 2--4 was chosen to capture both the core
two-character substitution pairs (\textit{ea}/\textit{ya},
\textit{oa}/\textit{wa}) and the longer consonant cluster
divergences (\textit{tjh}/\textit{ch}, \textit{kgh}/\textit{kh}).
Unigrams provide no discriminative signal since individual
characters are shared across both variants. Extending beyond
4-grams would increase feature dimensionality without adding
signal for the specific substitutions documented by
\citet{makutoane2022people}.

\paragraph{Diacritic preservation.}
The \texttt{strip\_accents} parameter was explicitly set to
\texttt{None}, overriding the scikit-learn default which strips
diacritical marks during vectorisation. This override is
linguistically motivated: LS orthography preserves diacritics
that SAS orthography omits \citep{matlosa2017sesotho}. Stripping
diacritics during vectorisation would destroy the orthographic
signal they carry, degrading classifier performance on the LS
class and producing a feature space that does not reflect the
actual textual properties of the two variants.

\paragraph{Class weighting and regularisation.}
The \texttt{class\_weight='balanced'} parameter adjusts sample
weights inversely proportional to class frequencies, compensating
for the 200 SAS / 102 LS imbalance in the training corpus without
requiring oversampling. Oversampling was avoided because it would
introduce artificial regularity into the LS character distribution
by replicating existing sentences, potentially causing the
classifier to overfit to specific headline constructions rather
than general orthographic patterns. The regularisation parameter
$C$ was set to 1.0, the scikit-learn default, which was retained
as no systematic overfitting was observed on the held-out test
set.

\paragraph{Evaluation metric.}
Macro-averaged F1 was used as the primary evaluation metric,
assigning equal weight to both classes regardless of their
frequency in the test set. This choice aligns with the pipeline's
design goal of handling both SAS and LS reliably, not merely
achieving high overall accuracy by exploiting the majority class.


\subsection{Data Processing and Sampling }

The verified and cleaned corpus was split into training and test
sets using an 80/20 ratio with stratified sampling, preserving the
SAS/LS class proportions in both partitions. Stratification was
applied because the 2:1 class imbalance would otherwise risk
concentrating LS sentences disproportionately in one split,
producing evaluation results that do not reflect the classifier's
generalisation across both variants. The split yielded 241 training
sentences (160 SAS, 81 LS) and 61 test sentences (40 SAS, 21 LS).
A fixed random seed (\texttt{random\_state=42}) was applied to
ensure reproducibility.

\subsection{Parameter Tuning}

No hyperparameter search was conducted beyond the design decisions
documented in Section~\ref{sec:learned}. The n-gram range, maximum
feature count (1,000), classifier type, and regularisation
parameter were set on the basis of the linguistic properties of
the task and established practices in the similar language
identification literature \citep{cianflone2016n,
de2024data} rather than through data-driven optimisation.
This approach was adopted to avoid the risk of overfitting
hyperparameters to a small held-out test set of 61 sentences, which
would provide unreliable estimates of generalisation performance.

\subsection{Evaluation Pipeline}
The classifier was trained by calling \texttt{pipeline.fit(X\_train,
y\_train)} on the 241-sentence training set, where the vectoriser
and classifier are fitted jointly within the scikit-learn
\texttt{Pipeline} object. Evaluation was performed by calling
\texttt{pipeline.predict(X\_test)} on the 61-sentence held-out test
set, with performance measured using macro-averaged F1, per-class
precision and recall, and a confusion matrix.

The rule-based baseline was evaluated on the same 61-sentence test
set using the marker-counting logic described in
Section~\ref{sec:rulebased}, enabling a direct comparison on
identical data. The interpretability of the learned classifier was
assessed by extracting the \texttt{LogisticRegression} coefficient
vector and identifying the character n-gram features with the
largest positive coefficients for each class, following the
interpretability-oriented evaluation approach adopted by
\citet{de2024data} in their LID model analysis.

\section{Summary}

\bibliographystyle{ruauthordate}
\bibliography{methodology}  
\end{document}
